\section{Introdução}

[DIRECIONAMENTO] Contextualizar a importância das APIs REST e o desafio de validar robustez além do happy path. Explicar que o TCC do aluno (Gomes, 2026) já implementa um fuzzer guiado por contrato OpenAPI, mas seu analisador atual classifica erros de forma binária (status code). Apresentar a lacuna: falta de um oráculo capaz de capturar anomalias sutis — como um 2xx com corpo inconsistente ou um 4xx com latência anormal. Formular a pergunta de pesquisa: "Como a lógica fuzzy pode melhorar a precisão na classificação de anomalias em relação a abordagens puramente booleanas?". Definir o objetivo geral (implementar e avaliar um FIS) e os específicos (definir variáveis, criar regras, validar com dados reais). Justificar a escolha da lógica fuzzy: a anomalia não é binária, há graus de criticidade, e a fuzzy modela essa nebulosidade naturalmente. Delimitar o escopo: o trabalho foca no módulo analisador (inteligência), tratando o fuzzer como provedor de dados. Finalizar com uma frase sobre a estrutura do artigo.
\endinput
